\section{Lipid và chất béo}

\subsection{Khái niệm lipid và chất béo}

Lipid là nhóm hợp chất hữu cơ có trong tế bào sống, thường không tan trong nước nhưng tan trong một số dung môi hữu cơ. Chất béo là một nhóm lipid quan trọng, có nhiều trong mỡ động vật và dầu thực vật.

Chất béo đơn giản là triester tạo bởi glycerol và các acid béo. Công thức tổng quát:
\[
\ce{(RCOO)3C3H5}
\]
Trong đó, $R$ là gốc hydrocarbon của acid béo. Một phân tử chất béo được tạo thành từ một phân tử glycerol và ba phân tử acid béo.

\begin{ghinho}{Cần phân biệt}
\begin{itemize}
\item \textbf{Lipid}: tên gọi chung của một nhóm hợp chất hữu cơ có tính chất vật lí đặc trưng.
\item \textbf{Chất béo}: lipid có cấu tạo là triester của glycerol và acid béo.
\end{itemize}
Không phải mọi lipid đều là chất béo.
\end{ghinho}

\subsection{Trạng thái tự nhiên và tính chất vật lí}

Chất béo có trong mô mỡ động vật, hạt và quả của nhiều loài thực vật. Ở điều kiện thường, một số chất béo ở thể rắn hoặc mềm như mỡ, bơ; một số khác ở thể lỏng như dầu ăn.

Chất béo:
\begin{itemize}
\item nhẹ hơn nước và không tan trong nước;
\item tan trong một số dung môi hữu cơ;
\item có thể tồn tại ở thể rắn hoặc lỏng tùy thành phần acid béo.
\end{itemize}

\subsection{Vai trò và cách sử dụng chất béo}

Lipid tham gia cấu tạo tế bào và có vai trò tích lũy năng lượng trong cơ thể. Chất béo còn được sử dụng trong chế biến thực phẩm, sản xuất xà phòng, glycerol và một số sản phẩm công nghiệp.

Chất béo cần thiết cho cơ thể nhưng nên sử dụng với lượng phù hợp. Chế độ ăn cần đa dạng thực phẩm, hạn chế ăn quá nhiều thức ăn giàu chất béo và kết hợp vận động hợp lí để bảo vệ sức khỏe.

\begin{vidu}
Một mẫu dầu ăn và một mẫu dầu bôi trơn đều không tan trong nước. Chỉ dựa vào tính tan trong nước có phân biệt được hai mẫu này không? Nêu một cách phân biệt dựa trên tính chất hóa học.
\end{vidu}
\begin{loigiai}
Chỉ dựa vào tính tan trong nước thì chưa phân biệt được vì cả hai mẫu đều không tan trong nước.

Có thể đun nóng từng mẫu với dung dịch \ce{NaOH}. Chất béo trong dầu ăn bị thủy phân trong môi trường kiềm, tạo glycerol và muối sodium của acid béo; dầu bôi trơn có bản chất chủ yếu là hydrocarbon nên không xảy ra phản ứng xà phòng hóa.

\textbf{Vậy,} có thể phân biệt hai mẫu bằng phản ứng với dung dịch \ce{NaOH} khi đun nóng.
\end{loigiai}
